\begin{center}
[2.]\par\medskip
\textsc{Gauss} to \textsc{Bessel}.\quad Göttingen, 26~December 1821.
\end{center}

\par\noindent\hspace*{\parindent}.\ .\ .\ .\ .\ .\quad Last summer was spent mostly on preparations for triangulation and on triangulation itself.\quad.\ .\ .\ .\quad The angles at the observatory, the meridian mark, the Hohehagen, the Hils and the Brocken have been measured. Only at the last point did the measurement of the angle between the \textls[120]{Inselsberg} and other points fail, because the extraordinarily bad weather while the heliotrope was there made almost all observation impossible.
The direction could be measured only once, for half an hour and with poor results, which do not agree well with other earlier and equally poor measurements, in which the instrument was aimed once at the house and once at \textsc{Encke}'s sextant auxiliary heliotrope 3 minutes before sunset.

\par\noindent\hspace*{\parindent}Having only one heliotrope at my disposal held things up enormously.
While I was on the Hils, the heliotrope had to be moved from Lichtenberg successively to the meridian mark, the Brelingerberg and the Deister. Similarly, while I was on the Brocken it went from Lichtenberg to the Hils, the Hohehagen and the Inselsberg.
Thus very few of the angles could be measured directly. At these two
\par\clearpage
\noindent350\hfill{\small\textbf{CORRESPONDENCE WITH BESSEL.}}\par\medskip
\noindent stations I had selected several other points that were visible in their own right, and I always measured whatever could be measured at the time.
I see from volume 4 of the Base du système métrique that the French did the same in Spain, but they did not combine the measurements correctly.
When I resume my triangulation next year, it will go better because I shall have two more heliotropes.
I always carried my sextant, which was fitted out as an auxiliary heliotrope, with me and used it to send telegraphic signals.
It works excellently, except that the auxiliary heliotrope is limited to angular separations below about \(138^{0}\), so I could not always use it at every time of day. Next year I would as a rule keep one of the heliotropes with me for this purpose.

\par\noindent\hspace*{\parindent}The greatest distance at which the heliotrope light has been seen with the naked eye is from the Brocken to the Hohehagen, where the telegraphic signals I sent from the Brocken were seen in this way. The distance is \(9\frac{1}{3}\) geographical miles, but favourable circumstances are probably required.
At distances of up to 6 miles, I saw the heliotrope light easily with the naked eye whenever circumstances were even moderately favourable (despite the loss of light when using the lorgnette). I likewise saw the heliotrope light from the Hils to the Brocken (\(7\frac{1}{2}\) miles) in the morning hours.
The mirror area is \(2\frac{1}{2}\) square inches. For the two new heliotropes I have used \(6\frac{1}{4}\) square inches. The light can easily be reduced, but at great distances there may be occasions when the stronger light is welcome. On one day, for example, the light from the Brelingerberg to the Hils got through only briefly, \textls[120]{although} (or, more accurately, \textls[120]{because}) the sky was completely cloudless.
On such days the atmosphere is as a rule least transparent. I remember one such day when the heliotrope light from the Deister was quite easy to see but could not be used, because none of the other objects could be seen, not even Einbeck, just over a mile away.

\par\noindent\hspace*{\parindent}Heliotrope light has a very substantial advantage over every other signal, and I have often made interesting observations on this point. The greater the magnification, the better one sees the heliotrope light, but the worse one sees terrestrial signals (at great distances), because with the latter it is chiefly their \textls[120]{faintness} that prevents them from being seen. Of Hanover I have,
\par\clearpage
\noindent{\small\textbf{ON THE HANOVER TRIANGULATION.}}\hfill351\par\medskip
\noindent for example, never been able to see a trace from the Brocken, even though I knew its exact position in elevation and azimuth.
On the day I received the faint heliotrope light from the Inselsberg (\textls[120]{solely} the fault of the weather, for with reasonably favourable air the light of a heliotrope with a mirror area of \(2\frac{1}{2}\) square inches would still have been more than bright enough), I could still see it quite well through the telescope of the theodolite, but could discern no trace of the mountain's outline through it. My assistant, on the other hand, could see that outline quite tolerably in a very poor telescope with low magnification, but could only rarely see the heliotrope light through it.
This is easily explained. Even at the greatest magnification the heliotrope light remains a point and its light is always the same, but the greater the magnification, the darker the background, etc.

\par\noindent\hspace*{\parindent}I built only two large signals, on the Hohehagen and the Hils, and I would probably have dispensed with them too if my heliotrope had been completed six weeks earlier.
From the Brocken, the Hils signal is seen against a nearby dark background. In the whole month I spent on the Brocken, I managed to \textls[120]{see} it at all only twice, and saw it well enough to have measured an angle for only about two minutes. The Hohehagen signal is seen against the more distant hills near Cassel. It was visible more often and could be observed a couple of times, although if I had not sent the heliotrope there I still would not have obtained all my angle measurements.
Apart from poor visibility, expense and the time they take, the signal towers have another very unfortunate drawback. They tempt people to acts of wanton destruction.
Unfortunately, my Hohehagen signal was recently almost completely vandalised, and I shall be glad if I can just find the point again with sufficient precision.

